\section{The classical high-probability birthday threshold}
\label{sec:classical-birthday}

In this section we prove that at \(A=C_q\) the classical high-probability threshold has the sharp prefactor \(\kappa_1=\sigma/(2\sqrt{2\pi})\) at scale \(\sqrt n\exp(C_q n^{1/3})\), with equality belonging to the failure regime. We treat three target types: ordinary vertices, outer vertices, and ordinary edges.

The strategy for ordinary vertices is a second-moment argument on a large family of vertex pairs at nearly typical distance. We first show that for typical pairs of candidate pairs, the joint collision probability factorizes asymptotically via Fourier localization and a Gaussian covariance comparison. We then bound the contribution of atypical pairs of pairs using a descending-radius argument, verify that the expected number of collisions diverges even at the critical prefactor, and finally adapt the proof to outer vertices and ordinary edges.

\step{Candidate pairs and their joint probability.}
We begin by defining our family of candidate pairs and the subset of typical configurations, with the goal of proving that any two typical pairs collide almost independently. Write \(s=n^{1/3}\) and \(K=pq^n\), and consider unordered vertex pairs at distance \(h\) nearest to the mean distance \(\mu=(1-1/q)n\), choosing \(h\) even when \(q=2\). The number of such candidate pairs is \(M=\exp(2Ln-O(\log n))\), and by symmetry they share a common ordinary collision probability \(P\).

We say that two candidate pairs \(x,y\) and \(u,v\) form a typical configuration if every ordered quadruple of letters from the alphabet appears among the coordinate tuples \((x_i,y_i,u_i,v_i)\) with frequency \(n/q^4+O(n^{11/12})\), for a fixed tolerance multiplier. Conditioning on the two within-pair distances being equal to \(h\) costs only a polynomial factor in probability, so standard concentration implies that the fraction of atypical configurations is \(\exp(-\Omega(n^{5/6}))\). Whenever a configuration is typical, all six pairwise distances among the four endpoints lie at \(\mu+O_q(n^{11/12})\).

\step{Disjoint singleton columns.}
To localize the characteristic function of the joint collision event, we first remove the linear dependencies among the shell discrepancy rows and construct disjoint families of words that perturb a single row at a time. In each of the two blocks (corresponding to the pairs \(x,y\) and \(u,v\)), the discrepancy rows sum to zero, so we delete one central discrepancy row per block. In the binary case with \(h\) even, the row sums vanish separately within each parity class, so we delete one central row of each parity instead. In the Fourier inversion integral, we set the Fourier frequency corresponding to each deleted row to zero.

We claim that within a central window of width \(O(n^{2/3})\) around \(\mu\), every remaining radius \(r\) admits disjoint sets of words that produce singleton columns in the two blocks, each with Poisson intensity at least \(\lambda_r\exp(-O(n^{1/6}+\log n))\). To construct the family for the first block, we prescribe words \(z\) satisfying
\[
 d(x,z)=r,\quad d(y,z)=c,\quad d(u,z)=d(v,z),
\]
where \(c\) is the deleted central row of the first block (of matching parity in the binary case).

To lower-bound the number of such words, consider the product law on words in which each coordinate independently disagrees with \(x\) with probability \(r/n\), choosing uniformly among the nonmatching letters. Conditional on the total distance from \(x\) being \(r\), this tilted product law is uniform on the Hamming sphere of radius \(r\) around \(x\). Partition the \(n\) coordinates into finitely many cells according to the coordinate tuple of the four endpoints, and round the expected symbol counts in each cell to integers while preserving the cell sizes. Under this tilted law, the conditional mean distance from \(z\) to any center at distance \(d\) from \(x\) is \(\mu+(1-qd/((q-1)n))(r-\mu)\). Because all six pairwise distances are \(\mu+O_q(n^{11/12})\) and \(r-\mu\) is at most \(O(n^{2/3})\), the rounded allocations satisfy the three constraints up to residuals of size \(O(1),O(1),O(n^{7/12})\), respectively.

We now adjust the symbol counts to eliminate these residuals. When \(q\ge3\), coordinates with endpoint tuples \((0,0,0,0),(0,1,0,0),(0,0,1,0)\) allow single-coordinate moves that change the triple of quantities \((d_x,d_y,d_u-d_v)\) by the vectors \((1,1,0),(0,1,0),(0,0,1)\), respectively, by changing the landmark letters \(0\to1,1\to2,1\to2\) in those cells. In the binary case, we translate coordinates so that \(x\) is the all-zero word and express the three constraints as linear equations on the counts of ones, with coefficient columns \((1,y_i,v_i-u_i)\) and right-hand sides \((r,(r+h-c)/2,(|v|-|u|)/2)\), all of which are integers. The cells with \((y_i,u_i,v_i)=(0,0,0),(1,0,0),(0,0,1)\) form a unimodular minor.

In both cases, we can cancel the residuals by making a bounded number of prescribed adjustments of size \(O(n^{7/12})\) inside cells of linear size, so every symbol count stays strictly interior to its cell. By Stirling's formula, the multinomial probability of the adjusted allocation under the tilted product law is at least \(\exp(-O(\log n+(n^{7/12})^2/n))\), and conditioning on the total distance \(r\) only increases this probability. Multiplying by the shell intensity \(\lambda_r\) gives the claimed Poisson intensity lower bound. Exchanging the roles of the two blocks produces the corresponding family for the second block, and by construction the singleton events for distinct radii and blocks are mutually disjoint.

\step{Joint intensity bound.}
Next we bound the Poisson intensity \(J_{r,t}\) of words lying simultaneously at radius \(r\) from one center and radius \(t\) from another distinct center in a typical configuration. For any two distinct typical centers and any \(r,t\) in the central window, we have the uniform bound
\[
 J_{r,t}\le \lambda_r\lambda_t K^{-1}\exp(o(s)).
\]
To prove this, apply the exponential Markov inequality to the exact probability of the atom \((r,t)\) using the signed tilts \((r-\mu)/(n\sigma^2),(t-\mu)/(n\sigma^2)\). Because the two centers are at distance \(\mu+O_q(n^{11/12})\) and both radii lie within \(O(n^{2/3})\) of \(\mu\), the mixed quadratic term in the cumulant generating function has size \(O(n^{1/4})=o(s)\), while the cubic remainder is \(O(1)\). By Stirling's approximation, the marginal shell exponents match the corresponding quadratic exponents up to \(O(\log n)\), which yields the stated upper bound.

\step{Fourier localization and Gaussian reduction.}
We now pass from the Bernoulli model to the Poisson approximation and localize the joint characteristic function to a small box on the dense rows, where the integrand simplifies to a product of small-row factors and a multivariate Gaussian. Replacing the Bernoulli log characteristic function by its Poisson counterpart introduces a uniform error of \(O(q^np^2)\), and discarding the tail frequencies outside the central window adds a further uniform error of \(o(1)\) to the log characteristic function.

Call a row dense if its shell intensity satisfies \(\lambda_r\ge e^{n^{1/4}}\), and set the box parameter \(H=e^{n^{1/5}}\). Whenever a dense coordinate lies outside the box \(|\theta_r|\le H/\sqrt{\lambda_r}\), the presence of a singleton column for that row bounds the modulus of the characteristic function by \(\exp(-\exp(\Omega(n^{1/5})))\). The same tail bound applies to the single-block Fourier integral, so in both integrals the discarded regions are negligible compared to the lower bound \(P^2\ge\exp(-2an-O(n^{2/3}\log n))\).

Inside the box, group the phase variables into three sets: the small rows of the first block, the small rows of the second block, and the combined dense rows of both blocks. When we split the Poisson exponent into three separate sums corresponding to these groups, the error is bounded by sums of products of individual phases drawn from different groups. Every such product involves distinct endpoints, so the joint intensity bound proved above controls its contribution. Because the \(\lambda_r\)-weighted sums of the absolute small phases and dense phases are at most \(n\exp(o(s))\) and \(O(nH\sqrt K)\), respectively, the total splitting error is \(o(1)\). Expanding the combined dense exponent in Taylor series, the linear term cancels by symmetry and the cubic remainder is \(O(nH^3e^{-n^{1/4}/2})\).

Consequently, on the box the joint integrand is uniformly \(1+o(1)\) times the product of the two single-block small-row characteristic functions and the dense Gaussian \(\exp(-\theta^TQ\theta/2)\), where \(Q\) denotes the dense Poisson covariance matrix. Each small-row factor is real and positive by the involution that swaps the two endpoints of its block. Having a positive proxy justifies converting uniform pointwise relative errors into relative errors for the integral, even though the joint integrand itself need not be real. Moreover, the disjoint singleton columns imply the lower matrix bound \(Q\succeq\exp(-O(n^{1/6}+\log n))\operatorname{diag}(\lambda_r)\), which allows us to extend the Gaussian integral from the box to all of Euclidean space via a union bound on Gaussian tails. By the same analysis, each single-block integral decomposes into the same small-row factor and the marginal Gaussian integral with covariance matrix \(Q_{ii}\).

\step{Radial basis and cross-Gram estimate.}
To show that the joint Gaussian integral factorizes into the product of the two marginal Gaussian integrals, we analyze the covariance matrices in a radial orthonormal basis and bound the cross-block correlation. In the uniform \(L^2\) space on a single coordinate, normalize the centered indicator of a given letter to have unit norm; the inner product between the normalized indicators of two different letters is \(-1/(q-1)\). For any center and each degree \(k\), define the unit degree-\(k\) radial vector as the normalized sum of coordinatewise products of these indicators over all \(k\)-subsets of coordinates. Vectors of distinct degrees are orthogonal, so these \(n+1\) vectors form an orthonormal radial basis around the center.

For two centers at distance \(d\), the inner product \(f_k(d)\) between their degree-\(k\) radial vectors equals the average over all \(k\)-subsets of the product of \(k\) entries drawn without replacement from a multiset containing \(d\) copies of \(-1/(q-1)\) and \(n-d\) copies of \(1\). For any fixed \(k\ge1\), comparing sampling without replacement to sampling with replacement incurs an error of \(O(k^2/n)\), so at typical distances \(f_k(d)\) tends uniformly to zero.

To control larger degrees, we separate the cases \(q\ge3\) and \(q=2\). When \(q\ge3\), applying the elementary symmetric mean inequality to the absolute values of the entries gives \(|f_k(d)|\le [1-(1-1/(q-1))d/n]^k\). When \(q=2\), the symmetry \(f_{n-k}(d)=(-1)^df_k(d)\) allows us to restrict to degrees up to \(n/2\). For \(k\le C\log n\), we use the sampling-with-replacement estimate. For \(C\log n\le k\le n/2\), we pair opposite signs into at least \(n/3\) disjoint pairs; only subsets that meet each such pair in zero or two coordinates can avoid cancellation in the average. Under independent Bernoulli sampling with parameter \(k/n\), the probability of avoiding all single-element intersections is at most \(\exp(-k/3)\), and conditioning on the subset having exact size \(k\) costs a factor of \(O(\sqrt k)\). Combining these bounds, we conclude that in both cases the sum of squared inner products is \(o(1)\), omitting degree 0 and, in the binary case, degree \(n\) (which vanishes identically on signed differences of even-distance pairs).

Now return to the covariance matrices. In each block, the signed shell difference vectors lie in the linear span of the differences between the radial vectors of the two endpoints. For every nonzero degree, the squared norm of this difference vector is \(2(1-f_k(h))\), which is bounded away from zero. Orthonormalizing the full block spans, the cross-Gram matrix between the two blocks is diagonal in the degree index and has Frobenius norm \(o(1)\), by the four cross-distance bounds on \(f_k\) just established. Restricting to the actual dense row subspaces corresponds to compressing this cross-Gram matrix by isometries, which cannot increase its Frobenius norm. Hence the whitened cross-covariance matrix \(T=Q_{11}^{-1/2}Q_{12}Q_{22}^{-1/2}\) has Frobenius norm \(o(1)\), and the ratio of the joint Gaussian integral to the product of the two marginal Gaussian integrals is \(\det(I-T^*T)^{-1/2}=1+o(1)\). Since the small-row factors cancel between the numerator and denominator, we obtain \(\Pr(\text{both collide})=(1+o(1))P^2\) uniformly over all typical configurations.

\step{Atypical pairs and exact cutoff.}
We now control the contribution of atypical configurations to the second moment, showing that they are negligible compared to \((MP)^2\), and then verify that \(MP\to\infty\) even at the exact critical prefactor \(\kappa=\kappa_1\).

First consider atypical pairs of candidate pairs. Whenever there is at most one cross-pair match within distance \(\delta n\), the balanced-cell two-pair bound from earlier already suffices, because the \(\exp(-\Omega(n^{5/6}))\) fraction of atypical configurations overwhelms the \(\exp(O(n^{2/3}\log n))\) loss in that bound.

Next suppose there are two cross-pair matches within distance \(\delta n\), and at least one of them has nonzero quotient distance (for \(q\ge3\), this simply means nonzero ordinary distance). Choose the two endpoints of that match as the anchors, one in each block. By the triangle inequality, the other two endpoint comparisons for each anchor have quotient distance bounded away from zero. We modify the rounded balanced allocations by \(O_q(1)\) coordinates so as to force the close comparison strictly inward toward \(\mu\). Specifically, on the positions where the anchor and its close comparator differ, the difference between their distances to a landmark is the count of anchor letters minus the count of comparator letters, which under the tilted shell law has mean \(-qd(r-\mu)/((q-1)n)\). When the separation \(d\) is sufficiently large, both letter counts are \(\Omega_q(d)\), so swapping a bounded number of letters from one to the other gives the difference the strict inward sign, and a compensating bounded adjustment on agreeing positions restores the total radius \(r\). If instead \(d=O_q(1)\), we place only comparator letters on the differing positions when \(r>\mu\), and only anchor letters when \(r<\mu\). Because all modifications are bounded, omitting a fixed central strip wide enough that these adjustments cannot cross \(\mu\) or disrupt the two far comparisons ensures that every comparison moves inward. In the binary case, we complement the close comparator first if needed, which preserves absolute deviation from the mean.

With these modifications, the balanced tilted cell probabilities remain polynomially large, and the same two-block descending-radius argument gives a joint collision probability of \(\exp(-2an+O(n^{2/3}\log n))\). The number of pairs of pairs with two close cross-matches is at most \(\exp((2L+2b_q(\delta))n+O(\log n))\), which becomes exponentially small after dividing by \(M^2P^2\) once \(\delta\) is chosen sufficiently small. Finally, if both cross-matches have zero quotient distance, there are only \(O(M)\) such pairs of pairs, and each has joint probability bounded by \(P\). Putting all cases together, whenever \(MP\to\infty\) the second moment of the collision count is \((1+o(1))(MP)^2\), which yields concentration and failure with high probability.

\step{Divergence of the expected count at equality.}
Our earlier lower bounds already imply \(MP\to\infty\) whenever \(a<2L\), as well as at \(a=2L\) when \(\beta<1/2\) or when \(\beta=1/2\) and \(\kappa<\kappa_1\). It remains to prove that \(MP\to\infty\) also holds at equality, \(\kappa=\kappa_1\). We do this by computing the leading terms of the Fourier Jensen lower bound for \(\log P\) through order \(s=n^{1/3}\) and showing that the savings from deleting one central row exceed the total excess cost across all remaining rows.

Set \(v=\sigma^2\), \(p_0=1-1/q\), \(R=\sqrt{2vA}\), and rescale the radius by \(x=(r-\mu)/n^{2/3}\). On the retained window at equality, Stirling's expansion gives the log shell means
\[
 h_r=\log(4\pi\lambda_r)=s(A-x^2/(2v))-c_3x^3
       +s^{-1}(b x-c_4x^4)+O(s^{-2}),
 \quad c_3=\frac{2p_0-1}{6v^2},\quad
 c_4=\frac{1-3v}{12v^3},
\]
where \(b\) is the linear Stirling correction. We claim that the half-sum of the positive parts \((h_r)_+\) satisfies
\[
 an+sR^5\left(\frac{c_3^2v}{2}-\frac{c_4}{5}\right)+o(s)
 =an-sR^5\frac{1+2v}{360v^3}+o(s).
\]
To obtain this expansion, replace the discrete sum by the integral of its clipped smooth approximation; the trapezoidal rule error, including the two corner intervals where the argument crosses zero, is \(o(s)\). On the symmetric interval \([-R,R]\), the odd terms integrate to zero. Near each endpoint \(x=\pm R\), the shift of the zero-crossing caused by the \(-c_3x^3\) term contributes a triangular region of area \((c_3R^3)^2/(2s(R/v))\) before multiplying by the half-density of shells, while the \(O(s^{-2})\) shell error sums to \(O(1)\) across the window.

Now apply the single-block Fourier Jensen bound with the central frequency already fixed to zero. Each remaining factor in the product is the average of \(e^{-c_n\lambda_r\theta^2}\) over \([-\pi,\pi]\), where \(c_n=1+O(p)\) has negligible effect. Compare the negative logarithm of each factor with \((h_r)_+/2\). When \(h_r<0\), Jensen's inequality bounds the excess by \(\pi e^{h_r}/12\); when \(h_r\ge0\), the excess equals \(-\log\operatorname{erf}(\sqrt{\pi e^{h_r}/4})\). Near each boundary \(h_r=0\), consecutive values of \(h_r\) change by \((R/v+o(1))/s\). Summing these two excess bounds over all rows in the window therefore costs at most \((2v/R)sJ+o(s)\), where
\[
 J\le\frac{\pi}{12}+
     \frac{e^{-\pi/4}}{(\pi/4)(1-e^{-\pi/4})}<1.34 .
\]
For the positive part \(h_r\ge0\), this integral bound follows from \(\operatorname{erfc}(z)\le e^{-z^2}\) and the elementary inequality \(e^h\ge 1+h\). In the rescaled Riemann sum, both bounds are integrable and dominate the tails because of the linear slopes at the two boundaries. Outside the retained window, Jensen's inequality and the shell tails bound the total cost by \(o(1)\).

On the other hand, deleting one central row already saves \(As/2+o(s)\) in the negative log probability. This savings strictly exceeds the total excess bound \((2v/R)sJ+o(s)\), because \(AR/(4v)=3L/(4v)\ge3\log2>1.34\). It follows that \(\log(MP)\ge c s\) for some constant \(c>0\), so the second-moment argument applies at equality as well.

\step{Outer vertices and ordinary edges.}
We finish the section by extending the threshold proof to outer vertices and ordinary edges.

For outer vertices, any ordinary collision between two vertices forces their endpoint selectors to agree, because the radius-zero shell must match. The contributions of the two endpoints to the discrepancy cancel both when both endpoints are selected and when neither is selected. By independence of the endpoint selectors from the remaining landmarks, the ratio of outer to ordinary collision probabilities is \(P_{\rm out}/P=(1-p)^2/((1-p)^2+p^2)=1-o(1)\). Since the ordinary joint collision bounds dominate the outer joint collision bounds, the second-moment concentration carries over directly to outer vertices.

For ordinary edges, we restrict to the parallel family of edges whose differing coordinate is the first coordinate, with letters \(\{0,1\}\). For any landmark, the distances to the two endpoints of an edge separate according to the common offset \(b\in\{0,1\}\) contributed by the first coordinate. Offset zero occurs with probability \(2/q\) and supplies both the singleton column constructions and the balanced-cell bounds. In the Gaussian comparison, the cross-Gram bound holds in the direct sum of the two orthogonal offset fibers, which increases the squared Frobenius norm bound by at most a factor of two and remains valid under compression to the actual dense rows. The joint intensity bounds follow by conditioning on the offset and shifting each projected shell radius by at most one. In the binary case, the offset is always zero and both parity gauges continue to apply.

Compared with the \(n\)-coordinate vertex shell probability, the exact edge shell probability differs by the ratio \(1-(r/n-p_0)/p_0\). Taking its logarithm adds only an odd linear term of order \(s^{-1}\) to the expansion of \(h_r\) at equality, which integrates to zero on \([-R,R]\) and does not affect the leading calculation. Passing to the projected dimension changes \(M\) by at most a constant factor up to its polynomial prefactor. This establishes the failure converse for all three target types, and hence simultaneously by a union bound. The matching high-probability success assertions follow immediately from the first moment.
